\section{CONCLUSÃO}

A seção de Conclusão do Experimento deve sumarizar os achados, discutir se os objetivos foram alcançados e analisar as possíveis fontes de erro e limitações do experimento.

Com base nos resultados simulados na Tabela \ref{tab:t1}, os coeficientes de dilatação linear experimentais para latão, alumínio e aço (ferro) foram \SI{2.00e-5}{\degree C^{-1}}, \SI{2.20e-5}{\degree C^{-1}} e \SI{1.20e-5}{\degree C^{-1}} , respectivamente. Na simulação, observou-se uma correspondência perfeita com os valores tabelados, resultando em erros relativos nulos.

Em um cenário experimental real, desvios dos valores tabelados são comuns e esperados. As fontes de erro podem incluir:
\begin{itemize}
    \item Imprecisões nas leituras do relógio comparador, possivelmente devido a erros de paralaxe.
    \item Erros na medição do comprimento inicial $L_0$.
    \item Variações na temperatura ambiente e na temperatura de ebulição da água, que é influenciada pela pressão atmosférica local.
    \item Perda de calor para o ambiente durante o aquecimento do tubo, que pode fazer com que a temperatura média do tubo seja ligeiramente inferior à temperatura do vapor.
    \item Imprecisões na calibração dos instrumentos de medição.
\end{itemize}
Para aprimorar a precisão do experimento, algumas medidas poderiam ser implementadas:
\begin{itemize}
    \item Realizar múltiplas medições para cada material e calcular a média, o que ajuda a mitigar o impacto de erros aleatórios.
    \item Empregar instrumentos de medição com maior precisão e menor incerteza intrínseca.
    \item Isolar termicamente o tubo durante o processo de aquecimento para minimizar a perda de calor.
    \item Levar em consideração a pressão atmosférica local para determinar o ponto de ebulição exato da água.
\end{itemize}
Apesar das potenciais fontes de erro, este experimento proporciona uma compreensão fundamental da dilatação linear e da importância dos coeficientes de dilatação para diferentes materiais, validando o modelo teórico apresentado.
